\subsection*{비교 대상과 입력 조건}
\textbf{CPII 기업 비교기는 같은 산업·서비스·시장·분기의 기업별 부담을 동일한 기준분포에 놓고 비교하는 입력형 계산 도구.} 회사명 검색이나 외부자료 자동 수집 기능이 아니라, 사용자가 구성한 10개 입력변수의 정규화·합산·차이 분해 기능. v2의 변경 범위는 입력 검증과 결과 표시의 보강이며, v1의 활성 변수·방향·가중치·고정 기준분포는 유지.

\begin{itemize}[leftmargin=1.8em,itemsep=0.3em]
\item 관측단위: 기업-서비스-시장-분기. 비교집단의 식별 필드는 \nolinkurl{sector}, \nolinkurl{service_scope}, \nolinkurl{focal_market}, \nolinkurl{reporting_period}, \nolinkurl{benchmark_id}. 다섯 필드의 대소문자까지 일치하는 관측끼리의 비교.
\item 필수 입력: 규제 3개, 위험 4개, 시장 3개의 활성 변수 전부. 빈 셀은 실제 관측값 0과 구별하며, 불완전 입력의 순위·점수차 산출 차단. 과거 결과의 가용항목 평균과 달리 모든 비교기업에 동일한 10개 성분을 요구하는 완전입력 사양.
\item 입력 단위: 비율의 1\%는 \(0.01\), 배수의 3배는 \(3\). 세율·조달비용의 격차는 기업값에서 공통 기준값을 차감하고, IP 자산화 결손은 공통 무형자산/매출 기준에서 기업값을 차감하는 방향.
\item 출처 기록: 변수별 원천, 기준일, 산식, 직접값·가공값·대리값 구분. \nolinkurl{company_exposure_factor}는 시장공통값에 이미 적용한 기업별 노출조정의 기록이며, 계산기가 입력변수에 다시 곱하는 계수는 아님.
\item 등록 상태: 동봉된 \nolinkurl{BM-DIGITAL-CONTENT-2025Q4}는 합성 예시 전용 기준. 실제 기업 적용의 선행조건은 분야·분기별 실제 비교군, 격차 산식, 출처 해시의 등록과 계산파일 재생성. 합성 표식을 실제 기업에 붙여 입력 검증을 우회하는 사용의 배제.
\end{itemize}

\subsection*{고정 정규화와 점수차 분해}
기업 \(a\)의 변수 \(j\) 원값을 \(x_{j,a}\), 부담방향 부호를 \(d_j\), 변환함수를 \(T_j\)로 표시. 투자회수율과 세전이익 회수 대리값에 \(d_j=-1\), 나머지 활성 변수에 \(d_j=1\) 적용. 해외의존비율에는 항등변환, 나머지 9개에는 부호보존 로그변환 적용.
\[
T_j(z)=
\begin{cases}
z,&j=\mathrm{foreign},\\
\operatorname{sgn}(z)\log(1+|z|),&\text{그 밖의 활성 변수},
\end{cases}
\qquad
w_{j,a}=\min\{u_j,\max\{\ell_j,T_j(d_jx_{j,a})\}\}.
\]
\(\ell_j,u_j\)는 분야-국가-분기 결합 기준표본의 변환값에서 선형보간으로 산출한 1\%·99\% 분위수. 이 경계와 변수별 기준표본은 비교기업을 추가하거나 삭제해도 불변. \(n_j\)개의 고정 윈저화 기준값을 \(w^0_{j,r}\)로 표시할 때 적용하는 순위식은 다음과 같음.
\[
L_j(w)=\sum_{r=1}^{n_j}\mathbf1\{w^0_{j,r}<w\},\qquad
U_j(w)=\sum_{r=1}^{n_j}\mathbf1\{w^0_{j,r}\leq w\},
\]
\[
S_{j,a}=\frac{L_j(w_{j,a})+U_j(w_{j,a})-1}{2(n_j-1)},
\qquad
I^{k,\mathrm{cmp}}_a=\frac{100}{m_k}\sum_{j\in J_k}S_{j,a},
\qquad
C_a=\frac13\sum_k I^{k,\mathrm{cmp}}_a.
\]
동일 기준값에는 평균순위를 부여하고, 기준값 사이의 새 입력에는 위 계수식으로 위치를 부여하는 방식. 연속적인 비용률 변화가 언제나 연속적인 점수 변화로 이어지지는 않으며, 같은 순위 구간이나 윈저화 경계 밖의 서로 다른 원값에 같은 점수의 부여 가능성.

\begin{itemize}[leftmargin=1.8em,itemsep=0.3em]
\item \(j,a,b,r,z,w\): 변수, 두 비교기업, 기준관측의 인덱스, 변환 전 실수 인수, 윈저화 후 위치 인수.
\item \(x_{j,a},d_j,T_j,w_{j,a},\ell_j,u_j,w^0_{j,r},n_j\): 기업 원값, 방향 부호, 변환함수, 기업 윈저화 값, 하단·상단 경계, 고정 기준값, 변수별 기준관측 수.
\item \(L_j,U_j,\mathbf1\): 기준값 미만·이하의 누적 개수와 조건 성립 시 1인 지시함수. \(\operatorname{sgn}\)은 부호함수.
\item \(J_k,m_k,k\): 하위지수별 활성 변수 집합, 그 원소 수, 규제·위험·시장 하위지수 인덱스. \((m_{\mathrm{reg}},m_{\mathrm{risk}},m_{\mathrm{market}})=(3,4,3)\).
\item \(S_{j,a},I^{k,\mathrm{cmp}}_a,C_a\): 0--1 변수별 기준분포 위치, 0--100 비교기 하위지수, 0--100 비교기 CPII. 위첨자 \(\mathrm{cmp}\)는 비교기 사양의 표시.
\end{itemize}

\[
\Delta^{\mathrm{cmp}}_{ab}=C_a-C_b
=\sum_k\sum_{j\in J_k}c_{j,ab},
\qquad
c_{j,ab}=\frac{100}{3m_k}(S_{j,a}-S_{j,b}),\quad j\in J_k.
\]
\(\Delta^{\mathrm{cmp}}_{ab}\)는 A에서 B를 뺀 점수차, \(c_{j,ab}\)는 해당 변수의 점수 기여도. 규제·시장 변수 각각의 유효가중치는 \(1/9\), 위험 변수 각각은 \(1/12\). \textbf{양의 점수차는 A의 부담지수가 B보다 높다는 뜻이며, 개선의 뜻은 아님.} 반면 같은 기업의 개선량을 개선 전 점수에서 개선 후 점수를 뺀 값으로 정의하면 양수는 부담 감소의 뜻. 서로 다른 차감 순서의 구별.

\subsection*{합성 예시와 경제적 해석}
\begin{table}[H]
\centering
\small
\caption{비교기 합성 예시의 점수와 하위지수 기여도}
\begin{tabular}{lrrrr}
\toprule
지표 & 가온 콘텐츠(A) & 누리 미디어(B) & A--B & 총점차 기여 \\
\midrule
규제 하위지수 & 12.738 & 5.086 & 7.652 & 2.551 \\
위험 하위지수 & 73.358 & 65.062 & 8.296 & 2.765 \\
시장 하위지수 & 33.591 & 26.504 & 7.087 & 2.362 \\
CPII & 39.896 & 32.217 & 7.678 & 7.678 \\
\bottomrule
\end{tabular}
\par\smallskip
\begin{minipage}{0.96\linewidth}\footnotesize
실제 기업에 관한 관측·추정이 아닌 네 개 합성 예시 중 두 기업의 계산값. 하위지수 점수차의 3분의 1이 총점차 기여도이며, 반올림 전 값에 기초한 합계 검증.
\end{minipage}
\end{table}

이 예시의 \(7.678\)점은 세 경로 모두에서 A의 정규화 부담이 더 높다는 신호. 소비자 피해율 7.678\%, 위법 확률 7.678\%, 기업분포상의 백분위 차이와는 구별. 일반적으로 총점차가 양수여도 개별 변수 기여도는 음수일 수 있으므로, 총점 다음에 하위지수와 변수별 양·음 기여도의 확인 필요.

화폐금액은 점수차가 아니라 원단위 차이와 비교 가능한 경제적 노출액의 결합. 예를 들어 망비용률을 \(n_a,n_b\), 같은 기간·통화·가격기준의 매출액을 \(R_a,R_b\)로 놓으면 공통규모 비용 상당액은 다음과 같음.
\[
\bar R_{ab}=\frac{R_a+R_b}{2},\qquad
E^{\mathrm{net}}_{ab}=(n_a-n_b)\bar R_{ab}.
\]
\(\bar R_{ab}\)는 두 매출액의 산술평균, \(E^{\mathrm{net}}_{ab}\)는 공통 매출규모를 적용한 망비용 차이 상당액. 합성 예시의 비용률 \(0.180\)과 \(0.105\)에 공통 매출규모 1조원을 가정하면 \(0.075\times1\text{조원}=750\text{억원}\). 7.678 CPII점 전체의 금액환산이나 실제 손해액이 아닌, 한 채널의 조건부 회계 상당액.

인프라 결합값·해외의존도·비용/설비투자 비율의 별도 화폐환산은 비활성. 같은 비용을 공유하는 변수별 금액의 자동합산 금지. 조달비용의 기간환산은 입력 금리의 연율·분기율과 노출기간의 정합성을 전제로 하며, 이미 분기화된 비율에 분기계수를 다시 곱하는 중복 조정의 배제. 입력·노출액의 하한·기준·상한은 확률구간이 아닌 결정론적 민감도 가정.

\subsection*{v2 변경사항과 사용 순서}
\begin{itemize}[leftmargin=1.8em,itemsep=0.3em]
\item 인프라 관계식 검사 수정: v1의 열 순서 오인을 제거하고 변수명에 따라 망비용률·해외의존도·인프라 결합값의 위치를 식별. 점수 계산의 변경이 아니라 잘못된 관계식 경고의 수정.
\item 원본 예시 검증의 분리: 원래 관측 ID와 10개 값이 모두 같은 경우에만 고정 기대점수와 대조. 정상적인 예시 입력 변경은 새로운 계산으로 허용하고, 원본 예시 대조 대상에서 제외.
\item 입력조건의 명확화: 실제 달력에 존재하는 2000--2100년의 출처일, 영문 대문자·숫자·하이픈·밑줄·마침표로 구성된 1--64자 관측 ID, 지정된 분류값, 비교집단의 공백과 대소문자 검사. 임의의 한국어 회사명과 제한된 관측 ID의 구별.
\item 결과 보호: 입력 검증 실패 시 해당 기업의 비교화면 점수·짝비교 차이·기여도·금액·민감도 표시 차단. 비교집단이 다르면 전체 순위 대신 집단별 순위와 비교불가 표시.
\item 사용 순서: \nolinkurl{01_INPUT_COMPANIES}에 기업 입력, \nolinkurl{07_INPUT_PROVENANCE}에 변수별 근거 기록, 필요 시 \nolinkurl{02_INPUT_ECON_BASES}에 경제 노출액 입력, \nolinkurl{03_COMPARE_2}에서 두 기업의 점수·기여도 확인, \nolinkurl{04_COMPARE_N}에서 집단별 순위 확인, \nolinkurl{06_SENSITIVITY}와 \nolinkurl{99_QA_AUDIT}에서 민감도·검증 상태 확인.
\item 수용 한도: 기업 100개, 화면상의 짝비교 행렬 30개. 더 큰 집합의 비교 및 입력변수별 범위 민감도는 동봉 Python 엔진의 별도 실행. Excel 민감도 화면은 가중치·저지원 제외·중복채널 축소의 네 가지 사양 비교.
\end{itemize}

활성 10개 변수를 동시에 관측한 기준행은 전체 159행 중 2행이며, 시장간 전이와 해외의존의 저지원 변수 가중치는 합계 \(2/9=22.22\%\). 따라서 v2의 지위는 연구용 베타이며, 입력 검증 개선을 기준표본 확대나 새로운 인과적 식별로 해석하는 것은 부적절. 비용회수·투자회수·시장전이 경로를 분리하여 비교하고 반대 방향의 기여도까지 드러내는 것이 이 도구의 실질적 용도. 기존 실증 결과의 기업 평균, 한국시장 방향점수, 정책 시나리오 감소량을 합성 예시의 점수로 대체하는 사용의 배제.
